problem:  
Let \((\mathcal{L}^{n+1}, \bar{g})\) be a globally hyperbolic Lorentzian manifold admitting a smooth, future-directed, null vector field \(\bar{V}\) that generates a **shear-free, twist-free, expanding null congruence**, i.e. in the usual optical decomposition, the shear tensor vanishes, the rotation tensor vanishes, and the expansion scalar \(\theta = \operatorname{div}_{\bar{g}}(\bar{V})\) is positive.  
Let \(M^n \subset \mathcal{L}^{n+1}\) be a complete spacelike hypersurface, with induced Riemannian metric \(g\), Levi-Civita connection \(\nabla\), second fundamental form \(h\), and shape operator \(A\).  

Assume:
1. There exists a smooth function \(f : M \to \mathbb{R}\) such that \((M,g,f)\) is a **steady gradient Ricci soliton** satisfying  
   \[
   \operatorname{Ric}_g + \nabla^2 f = 0.
   \]
2. The tangential projection \(V = (\bar{V})^T\) of \(\bar{V}\) onto \(M\) coincides with the soliton vector field: \(V = \nabla f\).  

**Conjecture / Open Problem:**  
Under these hypotheses, must the intrinsic gradient field \(\nabla f\) be a **conformal vector field** on \((M,g)\)?  
Equivalently, does the ambient optical condition (shear-free, twist-free, expanding \(\bar{V}\)) imply that
\[
\nabla_i V_j + \nabla_j V_i = 2\phi\, g_{ij}
\]
for some smooth scalar field \(\phi\) related to the optical expansion \(\theta\) and the mean curvature \(H = \operatorname{tr}_g A\) of the hypersurface?

In more geometric terms: can the *shear-freeness* of the ambient null congruence enforce that the soliton potential flow lines on \(M\) are conformally self-similar, thus constraining the extrinsic mean curvature evolution via the Gauss–Codazzi and optical scalar equations?  

---

Why is it a "good" problem:  
This version isolates a single, sharply-defined conjecture: whether the **shear-free condition in Lorentzian optical geometry forces conformality** of the induced soliton potential on a spacelike hypersurface. The hypothesis involves well-understood geometric quantities—optical expansion, shear, and mean curvature—and connects them via projection relations derived from the Gauss–Codazzi system. The question probes a new potential mechanism by which **causal optical structures** in spacetime could produce **intrinsic conformal geometry** on Ricci soliton hypersurfaces.  

It lies squarely at the intersection of **Ricci flow theory**, **Lorentzian submanifold optics**, and **geometric analysis**—fields that rarely meet in a single conjecture. The problem is elegantly simple to state (“Does shear-free expansion imply conformality of the soliton vector field?”) yet would require the integration of analytic and geometric techniques to resolve. A positive (or negative) resolution would either uncover a new rigidity property uniting soliton dynamics and null congruence optics or identify the precise analytical obstruction that prevents such correspondence, opening a novel cross-disciplinary avenue between general relativity’s optical geometry and Riemannian Ricci soliton theory.